\documentclass[11pt,a4paper]{article}

\usepackage[utf8]{inputenc}
\usepackage[T1]{fontenc}
\usepackage[ngerman]{babel}
\usepackage[a4paper,left=22mm,right=22mm,top=18mm,bottom=20mm]{geometry}
\usepackage{xcolor}
\usepackage{hyperref}

% HENSOLDT-Farben und Schrift passend zum Lebenslauf.
\definecolor{CVPrimary}{HTML}{f2581d}
\definecolor{CVText}{HTML}{222222}

\hypersetup{colorlinks=true,urlcolor=CVPrimary}
\renewcommand{\familydefault}{\sfdefault}
\pagestyle{empty}
\setlength{\parindent}{0pt}
\setlength{\parskip}{0.65em}
\setlength{\emergencystretch}{2em}

% Für ein festes Versanddatum beispielsweise: 14. Oktober 2026
\newcommand{\Briefdatum}{\today}

\begin{document}
\color{CVText}

{\raggedleft \Briefdatum\par}
\vspace{5mm}
{\large\bfseries\textcolor{CVPrimary}{Data Analyst – Maschinendaten \& AI}}\par
{\color{CVPrimary}\rule{\linewidth}{0.6pt}}\par
\vspace{1mm}

Sehr geehrte Damen und Herren, 

Aus Daten werden bessere Produkte und neue Services: Genau diesen Schritt reizt mich an der Stelle als Data Analyst bei Orbitalum. Maschinen-, Prozess- und Servicedaten systematisch zu nutzen, um Servicebedarfe früh zu erkennen und neue datenbasierte Funktionen zu entwickeln, ist die Art von Arbeit, die ich machen möchte.

Bei ZEISS SMT habe ich als Werkstudentin Maschinendaten zur Verfügbarkeit analysiert und überwacht, Reports und Dashboards in Power BI gebaut und die Ergebnisse eng mit den Fachabteilungen abgestimmt. Dort habe ich erlebt, dass Analysen erst dann wirken, wenn sie zu den Fragen der Fachbereiche passen. Im Master Machine Learning \& Data Analytics an der Hochschule Aalen habe ich mit Python vor allem Modelle für Zeitreihen und Bilddaten trainiert und getunt. Mein duales Bachelorstudium der Wirtschaftsinformatik bei Bosch hat mir gezeigt, wie Daten und IT-Systeme in einem Unternehmen zusammenhängen, und mir geholfen, Datenmodelle und KPIs so zu gestalten, dass sie für Produktmanagement und Service nutzbar sind.

Aktuell baue ich als wissenschaftlicher Mitarbeiter in einem Forschungsprojekt die Erfassung von Vibrationsdaten an einer Bandsäge auf. Sie soll später die Grundlage für Predictive Maintenance bilden.

Ich möchte gemeinsam mit Entwicklung, Produktmanagement und Service aus Maschinendaten konkrete Erkenntnisse ableiten und die Data-Analytics- und AI-Strategie bei Orbitalum mitgestalten. Über ein persönliches Gespräch freue ich mich.

Mit freundlichen Grüßen

\vspace{6mm}
Natalie Böhm

\end{document}
